% Copyright 2026 The terCAD team. All rights reserved.
% Use of this content is governed by a CC BY-NC-ND 4.0 license that can be found in the LICENSE file.

\subsubsection{Apple Store - macOS} \label{d-macos}

The process of creating \q{macOS}-package for a distribution is almost the same as for \q{iOS} (\ref{d-ios}) with a few
deviations. The desktop application is a Compose Multiplatform (JVM) one, so that it's packed by Gradle
(\q{./gradlew :desktopApp:packagePkg}) together with a Java runtime into \q{zwyka.app}, and then into a signed installer
\q{.pkg} (\q{macOS}-target of the \q{build}-matrix in \q{.github/workflows/build.yml}). The Mac App Store specifics are
declared in \q{desktopApp/build.gradle.kts}, and they are applied only when the provisioning profile is present (it's
added by the pipeline), so that local builds stay unaffected:

\begin{lstlisting}[language=kotlin]
macOS {
    bundleID = "com.tercad.zwyka"
    appCategory = "public.app-category.productivity"
    (findProperty("versionCode") as String?)?.let { packageBuildVersion = it }
    val appStoreProfile = project.file("embedded.provisionprofile")
    if (appStoreProfile.exists()) {
        appStore = true
        provisioningProfile.set(appStoreProfile)
        runtimeProvisioningProfile.set(project.file("runtime.provisionprofile"))
        entitlementsFile.set(project.file("entitlements.plist"))
        runtimeEntitlementsFile.set(project.file("runtime-entitlements.plist"))
    }
}
\end{lstlisting}

\noindent The main difference with \q{iOS} is that the application consists of two signed parts: the launcher
(\q{zwyka.app}) and the bundled Java runtime, so that both of them need their own identifier and provisioning profile.

\paragraph{Step 1: Certificates.} Unlike iOS, two certificates are required (the private key and the CSR-file from
\ref{d-ios} can be reused), both are taken from
\href{https://developer.apple.com/account/resources/certificates/add}{https://developer.apple.com/account/resources/certificates/add}:

\begin{itemize}
  \item "\emph{Mac App Distribution}" -- gives "3rd Party Mac Developer Application" certificate, used to sign the
  application and the runtime;
  \item "\emph{Mac Installer Distribution}" -- gives "3rd Party Mac Developer Installer" certificate, used to sign the
  \q{.pkg}-file.
\end{itemize}

\noindent Both of them are converted into \q{.p12} (with the same export password as for iOS) and into base64:

\begin{lstlisting}[language=terminal]
$ openssl x509 -in mac_app.cer -inform DER -out mac_app.pem -outform PEM
$ openssl pkcs12 -export -legacy -out mac_app.p12 -inkey apple.key -in mac_app.pem
$ base64 mac_app.p12 > apple-mac.txt
$ openssl x509 -in mac_installer.cer -inform DER -out mac_installer.pem -outform PEM
$ openssl pkcs12 -export -legacy -out mac_installer.p12 -inkey apple.key -in mac_installer.pem
$ base64 mac_installer.p12 > apple-mac-installer.txt
\end{lstlisting}

\noindent The signing identity for Gradle (\q{APPLE\_MAC\_IDENTITY}) is not issued separately: it's the common name
(\q{CN}) of the downloaded "Mac App Distribution" certificate without the "3rd Party Mac Developer Application: "
prefix. Compose adds the prefix by itself ("3rd Party Mac Developer Application: " to sign the application,
"3rd Party Mac Developer Installer: " to sign the \q{.pkg}), so that both certificates have to be issued for the same
account, which is the case. The name is read from the \q{.cer}-file:

\begin{lstlisting}[language=terminal]
$ openssl x509 -in mac_app.cer -inform DER -noout -subject
subject=UID=8PVKPQ758L, CN=3rd Party Mac Developer Application: terCAD (8PVKPQ758L), OU=8PVKPQ758L, O=terCAD, C=XX
\end{lstlisting}

\noindent Here \q{APPLE\_MAC\_IDENTITY} is "\emph{terCAD (8PVKPQ758L)}": the account holder name (the legal entity
name of an organization, or the full name of an individual, as shown in "Membership details" on
\href{https://developer.apple.com/account}{https://developer.apple.com/account}) followed by the Team ID in brackets.
On macOS the same value is listed by \q{security find-identity -v -p codesigning} after the \q{.p12} is imported.

\paragraph{Step 2: Identifiers.} The App ID \q{com.tercad.zwyka} from \ref{d-ios} is used for macOS as well (App IDs are
shared across Apple platforms). In addition, the Java runtime requires a second App ID. The reason is in the structure
of the package that is produced by \q{jpackage} (used by Compose): the JVM is a nested bundle
\q{zwyka.app/Contents/runtime} with its own \q{Info.plist}, and its identifier is the identifier of the application with
the \q{com.oracle.java.}-prefix. The Mac App Store requires each signed bundle inside the package to have a provisioning
profile that matches its identifier, so that the runtime gets its own App ID and its own profile.

\noindent The App ID is registered on
\href{https://developer.apple.com/account/resources/identifiers/list}{https://developer.apple.com/account/resources/identifiers/list}:

\begin{enumerate}
  \item "\emph{+}" (Register a new identifier) $\rightarrow$ "\emph{App IDs}" $\rightarrow$ "\emph{Continue}";
  \item type "\emph{App}" $\rightarrow$ "\emph{Continue}";
  \item Description: "Zwyka Runtime"; Bundle ID: "\emph{Explicit}" with the value
  \q{com.oracle.java.com.tercad.zwyka} (exactly the Bundle ID of the application, prefixed by \q{com.oracle.java.});
  \item no capabilities are selected $\rightarrow$ "\emph{Continue}" $\rightarrow$ "\emph{Register}".
\end{enumerate}

\paragraph{Step 3: Provisioning profiles.} Two profiles of the same type are created on
\href{https://developer.apple.com/account/resources/profiles/list}{https://developer.apple.com/account/resources/profiles/list},
first one for the application, second one for the runtime (\q{APPLE\_MAC\_RUNTIME\_PROFILE}):

\begin{enumerate}
  \item "\emph{+}" (Generate a profile) $\rightarrow$ in the "\emph{Distribution}"-section select "\emph{Mac App Store
  Connect}" $\rightarrow$ "\emph{Continue}";
  \item Profile Type: "\emph{Mac}" (not "Mac Catalyst"); App ID: \q{com.tercad.zwyka} for the first profile, and
  \q{com.oracle.java.com.tercad.zwyka} ("Zwyka Runtime" from Step 2) for the second one $\rightarrow$
  "\emph{Continue}" (there is no separate type for the runtime: it's the same "Mac App Store Connect" profile, and the
  App ID is listed only after it's registered in Step 2);
  \item certificate: the "\emph{Mac App Distribution}" one from Step 1 (shown as "3rd Party Mac Developer
  Application", not "Apple Distribution") $\rightarrow$ "\emph{Continue}";
  \item Provisioning Profile Name: "\emph{Zwyka macOS}" for the first profile, "\emph{Zwyka macOS Runtime}" for the
  second one $\rightarrow$ "\emph{Generate}" $\rightarrow$ "\emph{Download}".
\end{enumerate}

\noindent The downloaded files have the \q{.provisionprofile}-extension (a signed property list). The identifier inside
can be checked before uploading to GitHub (\q{security cms -D -i <file>} does the same on macOS):

\begin{lstlisting}[language=terminal]
$ openssl smime -inform der -verify -noverify \
    -in Zwyka_macOS_Runtime.provisionprofile 2>/dev/null \
    | grep -A1 application-identifier
<key>com.apple.application-identifier</key>
<string>8PVKPQ758L.com.oracle.java.com.tercad.zwyka</string>
\end{lstlisting}

\noindent Unlike iOS, the profiles are not installed on the runner, but embedded into the package by Gradle:
\q{embedded.provisionprofile} of \q{desktopApp/} goes to \q{zwyka.app/Contents/}, and \q{runtime.provisionprofile} --
to the runtime bundle. Both files are converted into base64 for the GitHub secrets:

\begin{lstlisting}[language=terminal]
$ base64 Zwyka_macOS.provisionprofile > apple-mac-profile.txt
$ base64 Zwyka_macOS_Runtime.provisionprofile > apple-mac-runtime.txt
\end{lstlisting}

\noindent Both profiles expire together with the "Mac App Distribution" certificate (in a year): then they have to be
regenerated (Edit $\rightarrow$ select the new certificate $\rightarrow$ Generate), and the secrets replaced.

\paragraph{Step 4: Application record.} The same application is extended by a new platform: on
\href{https://appstoreconnect.apple.com/apps}{https://appstoreconnect.apple.com/apps} open "Zwyka" and use
"\emph{Add Platform}" $\rightarrow$ "\emph{macOS}". Since the Bundle ID is the same, it's a universal purchase: the
name, privacy details and pricing are shared, while the macOS version has its own description, keywords, "What's New" and
screenshots (16:10 -- 1280x800, 1440x900, 2560x1600 or 2880x1800).

\paragraph{Step 5: GitHub secrets.} In addition to the secrets of \ref{d-ios} (\q{APPLE\_TEAM\_ID},
\q{APPLE\_P12\_KEY}, \q{APPLE\_API\_P8}, \q{APPLE\_API\_KEY}, \q{APPLE\_API\_ISSUE}) the following ones are required:

\begin{lstlisting}[language=terminal]
APPLE_MAC_P12              # content of apple-mac.txt (Step 1)
APPLE_MAC_INSTALLER        # content of apple-mac-installer.txt (Step 1)
APPLE_MAC_IDENTITY         # CN of the certificate without the prefix (Step 1)
APPLE_MAC_PROFILE          # content of apple-mac-profile.txt (Step 3)
APPLE_MAC_RUNTIME_PROFILE  # content of apple-mac-runtime.txt (Step 3)
\end{lstlisting}

\noindent After that, each pushed tag \q{v*} imports both certificates into a temporary keychain, puts the profiles into
\q{desktopApp/}, packs and signs \q{zwyka\_macOS.pkg}, attaches it to the release page and uploads it to App Store
Connect by \q{xcrun altool --upload-app --type macos} (the same step as for iOS):

\begin{lstlisting}[language=terminal]
- name: Run macOS Build
  if: matrix.target == 'macOS'
  run: |
    ./gradlew :desktopApp:packagePkg --no-daemon \
      -PversionName=${{ needs.release.outputs.version }} \
      -PversionCode=${{ needs.release.outputs.build_number }} \
      -Pcompose.desktop.mac.sign=true \
      -Pcompose.desktop.mac.signing.identity="${{ secrets.APPLE_MAC_IDENTITY }}" \
      -Pcompose.desktop.mac.signing.keychain="$MAC_KEYCHAIN"
\end{lstlisting}

\noindent \q{packagePkg} is used instead of \q{packageReleasePkg} to avoid ProGuard, the same as for other desktop
packages. The runner is Apple silicon, so the package is \q{arm64}-only: the application is available for Macs with
Apple silicon (M1 and later).

The version of the package is taken from the tag, and for macOS it has to start with 1 or higher: \q{jpackage} rejects
versions like \q{0.6.0} (tag \q{v0.6.0}), and \q{:desktopApp:createDistributable} fails with "External tool execution
failed". The reason is written only into the logs of \q{jpackage}
(\q{desktopApp/build/compose/logs/createDistributable/*-err.txt}):

\begin{lstlisting}[language=terminal]
Bundler Mac Application Image skipped because of a configuration problem:
  The first number in an app-version cannot be zero or negative.
Advice to fix: Set a compatible 'app-version' value.
  Valid versions are one to three integers separated by dots.
\end{lstlisting}

\paragraph{Sandbox.} Any application of the Mac App Store has to be sandboxed, otherwise the upload fails with "App
sandbox not enabled". Both parts declare \q{com.apple.security.app-sandbox}, while the runtime has to inherit the
sandbox of the application by \q{com.apple.security.inherit} (\q{desktopApp/runtime-entitlements.plist}). The JVM
requires JIT and unsigned executable memory, and Import / Export of data -- an access to the files that are chosen in the
file dialogs:

\begin{lstlisting}[language=xml]
<!-- ./desktopApp/entitlements.plist -->
<key>com.apple.security.app-sandbox</key>
<true/>
<key>com.apple.security.cs.allow-jit</key>
<true/>
<key>com.apple.security.cs.allow-unsigned-executable-memory</key>
<true/>
<key>com.apple.security.files.user-selected.read-write</key>
<true/>
<key>com.apple.application-identifier</key>
<string>TEAM_ID.com.tercad.zwyka</string>
\end{lstlisting}

\noindent \q{TEAM\_ID} is replaced by the pipeline from \q{APPLE\_TEAM\_ID}. The end-of-line character of the
entitlements has to be "LF", otherwise the list is not recognized.

The sandboxed application has no write access to the home folder, the same as Snap and Flatpak packages, so that the
database is placed into the container of the application (\q{\textasciitilde/Library/Containers/com.tercad.zwyka/Data}),
taken from \q{HOME} when \q{APP\_SANDBOX\_CONTAINER\_ID} is set (\q{desktopApp/src/main.kt}).

\paragraph{Step 6: Compliance and review.} \q{ITSAppUsesNonExemptEncryption} is declared as \q{false} by
\q{infoPlist} of \q{desktopApp/build.gradle.kts}, the same as for iOS. The build appears in the macOS section of the
"TestFlight"-tab; then it's selected on the macOS version page and submitted for review ("\emph{Add for Review}"
$\rightarrow$ "\emph{Submit for Review}"), the same as in \ref{d-ios}.
